\documentclass[11pt]{article}
\usepackage{mathblatt}
\begin{document}
% Kopfzeile Seite 1 ohne Kennung; die Rückseite (Lösungen, für den Lehrer) bekommt sie nach \begleitteil
\blattkopf{Mathematik}{Basisaufgaben}
\einheitenkopf[][Zettel 14]{Basisaufgaben · Zettel 14}
\zweigzeile{je Aufgabe 1 Punkt · Lösungen auf der Rückseite}
\par\vspace{2pt}\noindent Name \underline{\hspace{7cm}}\hfill Datum \underline{\hspace{3cm}}\hfill Punkte \underline{\hspace{1.2cm}} / 18\par\vspace{6pt}
% \small: volle Seite, kurze Aufgaben zweispaltig (v0.6)
\small

\noindent\begin{minipage}[t]{0.485\linewidth}\raggedright
% 1: Bruchteil einer Größe berechnen – brueche-dezimalzahlen-basis-k2-v5 (Original 2018-OS-B1a, ausgedacht: keine Marke)
\begin{aufgabe}{$\frac{3}{10}$ von $4$ kg – wie viel Kilogramm? \leerfeld[kg]}
\end{aufgabe}
\end{minipage}\hfill\begin{minipage}[t]{0.485\linewidth}\raggedright
% 2: Termwert berechnen – rationale-zahlen-basis-k1-v10 (Original 2026-FOR-B1g, ausgedacht: keine Marke)
\begin{aufgabe}{Berechne den Wert von $8 \cdot (x - 4)$ für $x = -1$. \leerfeld}
\end{aufgabe}
\end{minipage}\par

\noindent\begin{minipage}[t]{0.485\linewidth}\raggedright
% 3: Exponent einer Potenz bestimmen – potenzen-wurzeln-basis-k1-v5 (Original 2024-OS-B1g, ausgedacht: keine Marke)
\begin{aufgabe}{$10^x = 10\,000$ – welche Zahl ist $x$? x = \leerfeld}
\end{aufgabe}
\end{minipage}\hfill\begin{minipage}[t]{0.485\linewidth}\raggedright
% 4: Grundwert berechnen – prozentrechnung-basis-k1-v4 (Original 2025-OS-B1a, ausgedacht: keine Marke)
\begin{aufgabe}{Beim Kauf eines Spiels spart Can $7\,€$, das sind $50\,\%$ des alten Preises. Wie hoch war der alte Preis? \leerfeld[€]}
\end{aufgabe}
\end{minipage}\par

\noindent\begin{minipage}[t]{0.485\linewidth}\raggedright
% 5: Prozentwert berechnen – prozentrechnung-basis-k4-v10 (Original 2026-FOR-B1a, ausgedacht: keine Marke)
\begin{aufgabe}{$75\,\%$ von $48\,€$ \leerfeld[€]}
\end{aufgabe}
\end{minipage}\hfill\begin{minipage}[t]{0.485\linewidth}\raggedright
% 6: Zeiteinheiten umrechnen – einheiten-basis-k5-v2 (Original 2025-OS-B1b, ausgedacht: keine Marke)
\begin{aufgabe}{Ein Training dauert $0{,}2$ h. Gib die Dauer in Minuten an. \leerfeld[min]}
\end{aufgabe}
\end{minipage}\par

\noindent\begin{minipage}[t]{0.485\linewidth}\raggedright
% 7: Lineare Gleichung lösen – lineare-gleichungen-basis-k1-v8 (Original 2024-OS-B1d, ausgedacht: keine Marke)
\begin{aufgabe}{Löse die Gleichung $7 \cdot (x - 3) = -14$. x = \leerfeld}
\end{aufgabe}
\end{minipage}\hfill\begin{minipage}[t]{0.485\linewidth}\raggedright
% 8: Wahrscheinlichkeit einstufig – bruchrechnung-basis-k1-v5 (Original 2016-OS-B1f, ausgedacht: keine Marke)
\begin{aufgabe}{Ein Spielwürfel wird einmal geworfen. Wie groß ist die Wahrscheinlichkeit, keine 6 zu werfen? \leerfeld}
\end{aufgabe}
\end{minipage}\par

\noindent\begin{minipage}[t]{0.485\linewidth}\raggedright
% 9: Arithmetisches Mittel berechnen – daten-basis-k1-v3 (Original 2021-OS-B1f, ausgedacht: keine Marke)
\begin{aufgabe}{Lena hat diese Noten $2$; $3$; $1$; $2$; $2$. Berechne das arithmetische Mittel (Note). \leerfeld}
\end{aufgabe}
\end{minipage}\hfill\begin{minipage}[t]{0.485\linewidth}\raggedright
% 10: Median bestimmen – daten-basis-k3-v2 (Original 2024-OS-B1h, ausgedacht: keine Marke)
\begin{aufgabe}{Wartezeiten an der Kasse in Minuten: $12$; $7$; $9$; $15$; $7$; $10$. Bestimme den Median. \leerfeld[min]}
\end{aufgabe}
\end{minipage}\par

\noindent\begin{minipage}[t]{0.485\linewidth}\raggedright
% 11: Spannweite berechnen – daten-basis-k4-v1 (Original 2019-OS-B1i, ausgedacht: keine Marke)
\begin{aufgabe}{Eine Tafel Schokolade kostet in vier Läden: $1{,}29$; $0{,}99$; $1{,}59$; $1{,}19$ €. Gib die Spannweite an. \leerfeld[€]}
\end{aufgabe}
\end{minipage}\hfill\begin{minipage}[t]{0.485\linewidth}\raggedright
% 12: Zufallsgerät zu Wahrscheinlichkeit entwerfen – wahrscheinlichkeit-basis-k2-v2 (Original 2019-OS-B1g, ausgedacht: keine Marke)
\begin{aufgabe}{In einer Urne liegen $6$ rote Kugeln. Wie viele weiße Kugeln muss man dazulegen, damit $P(\text{rot}) = \frac{3}{4}$ ist? \leerfeld Kugeln}
\end{aufgabe}
\end{minipage}\par

% 13: Term zu Sachtext angeben – terme-basis-k2-v8 (Original 2023-OS-B1h, ausgedacht: keine Marke)
\begin{aufgabe}{Die Hälfte einer Zahl $x$ wird um 3 vergrößert. Welcher Term passt? \\ \kreuz{$(x + 3) : 2$} \kreuz{$x : 2 + 3$} \kreuz{$2x + 3$} \kreuz{$x : 3 + 2$}}
\end{aufgabe}

% 14: Scheitelpunktform aufstellen – quadratische-funktionen-basis-k4-v1 (Original 2016-OS-B1g, ausgedacht: keine Marke)
\begin{aufgabe}{Eine verschobene Normalparabel hat den Scheitelpunkt $S(2|-5)$. Welche Gleichung gehört zu ihr? Kreuze an.\\ \kreuz{$y = (x + 2)^2 - 5$}\\ \kreuz{$y = (x - 2)^2 + 5$}\\ \kreuz{$y = (x + 5)^2 + 2$}\\ \kreuz{$y = (x - 2)^2 - 5$}}
\end{aufgabe}

% 15: Rechteckseite aus Umfang berechnen – flaechen-basis-k1-v2 (Original 2018-OS-B1f, ausgedacht: keine Marke)
\begin{aufgabe}{Ein rechteckiges Foto hat einen Umfang von $24$ cm, eine Seite ist $7$ cm lang. Wie lang ist die andere Seite? \leerfeld[cm]}
\end{aufgabe}

% 16: Satz des Pythagoras formulieren – pythagoras-basis-k2-v1 (Original 2022-OS-B1g, ausgedacht: keine Marke)
\begin{aufgabe}{Im Dreieck $ABC$ liegt der rechte Winkel bei $C$. Welche Gleichung gilt? Kreuze an. \\ \kreuz{$a^2 = b^2 + c^2$} \kreuz{$a^2 + b^2 = c^2$} \kreuz{$a + b = c$}}
\end{aufgabe}

% 17: Graph einer linearen Funktion erkennen – lineare-funktionen-basis-k1-v2 (Original 2025-OS-B1f, ausgedacht: keine Marke)
\begin{aufgabe}{\begin{minipage}[t]{0.6\linewidth}Welcher Graph gehört zu einer linearen Funktion? Kreuze an. \\ \kreuz{A} \kreuz{B} \kreuz{C} \kreuz{D}\end{minipage}\hfill\begin{minipage}[t]{0.37\linewidth}\vspace{0pt}
\begin{ksys}[xmin=-5,xmax=5,ymin=-5,ymax=5,klein] \funktion{abs(\x+1)+1}{A} \funktionab[3.5]{-2/\x}{B}{0.4}{5} \parabel{-1}{0}{3}{C} \gerade{-1}{2}{D} \end{ksys}
\end{minipage}}
\end{aufgabe}

% 18: Term zu Figur angeben – flaechen-basis-k4-v2 (Original 2025-OS-B1i, ausgedacht: keine Marke)
\begin{aufgabe}{\begin{minipage}[t]{0.6\linewidth}Welche Formel beschreibt den Flächeninhalt $A$ des Quadrats? \\ \kreuz{$A = s^2$} \kreuz{$A = 2 \cdot s$} \kreuz{$A = s + s$} \kreuz{$A = 4 \cdot s$}\end{minipage}\hfill\begin{minipage}[t]{0.37\linewidth}\vspace{0pt}
\rechteck[punkte=,seiten={s,,,}]{2}{2}
\end{minipage}}
\end{aufgabe}

\begleitteil
\blattkopf{Basisaufgaben · Zettel 14}{BAS-Z14 · Lösungen}
\erg{1}{$4 : 10 \cdot 3 = 1{,}2$ kg \hfill {\footnotesize\textit{Bruchteil einer Größe berechnen · 2018}}}
\erg{2}{$8 \cdot (-1 - 4) = 8 \cdot (-5) = -40$ \hfill {\footnotesize\textit{Termwert berechnen · 2026}}}
\erg{3}{$10, 100, 1\,000, 10\,000$: vier Faktoren, also $x = 4$ \hfill {\footnotesize\textit{Exponent einer Potenz bestimmen · 2024}}}
\erg{4}{$7 \cdot 2 = 14\,€$ (nicht $50\,\%$ von $7\,€$) \hfill {\footnotesize\textit{Grundwert berechnen · 2025}}}
\erg{5}{$0{,}75 \cdot 48 = 36\,€$ \hfill {\footnotesize\textit{Prozentwert berechnen · 2026}}}
\erg{6}{$0{,}2 \cdot 60 = 12$ min \hfill {\footnotesize\textit{Zeiteinheiten umrechnen · 2025}}}
\erg{7}{$x - 3 = -2$, also $x = 1$ \hfill {\footnotesize\textit{Lineare Gleichung lösen · 2024}}}
\erg{8}{$\frac{5}{6}$ (die 1 bis 5) \hfill {\footnotesize\textit{Wahrscheinlichkeit einstufig · 2016}}}
\erg{9}{$10 : 5 = 2$ \hfill {\footnotesize\textit{Arithmetisches Mittel berechnen · 2021}}}
\erg{10}{geordnet $7$; $7$; $9$; $10$; $12$; $15$; $(9 + 10) : 2 = 9{,}5$ min \hfill {\footnotesize\textit{Median bestimmen · 2024}}}
\erg{11}{größter minus kleinster Wert: $1{,}59 - 0{,}99 = 0{,}60\,€$ \hfill {\footnotesize\textit{Spannweite berechnen · 2019}}}
\erg{12}{$6$ rote sind $\frac{3}{4}$ von $8$ Kugeln, also $8 - 6 = 2$ weiße \hfill {\footnotesize\textit{Zufallsgerät zu Wahrscheinlichkeit entwerfen · 2019}}}
\erg{13}{$x : 2 + 3$ \hfill {\footnotesize\textit{Term zu Sachtext angeben · 2023}}}
\erg{14}{$y = (x - 2)^2 - 5$ \hfill {\footnotesize\textit{Scheitelpunktform aufstellen · 2016}}}
\erg{15}{$24 : 2 = 12$; $12 - 7 = 5$ cm \hfill {\footnotesize\textit{Rechteckseite aus Umfang berechnen · 2018}}}
\erg{16}{$a^2 + b^2 = c^2$ \hfill {\footnotesize\textit{Satz des Pythagoras formulieren · 2022}}}
\erg{17}{D (eine Gerade ohne Knick) \hfill {\footnotesize\textit{Graph einer linearen Funktion erkennen · 2025}}}
\erg{18}{$A = s^2$ \hfill {\footnotesize\textit{Term zu Figur angeben · 2025}}}
\end{document}
